\section{Introduction}
\label{sec:introduction}
\pagenumbering{arabic}
\setcounter{page}{1}

Multimodal high-dimensional signals---including hyperspectral images, dynamic medical tomograms, multispectral video sequences, and photometric collections---form the foundation of modern Earth remote sensing, medical diagnostics, and computational vision \citep{bioucas2012hyperspectral, kolda2009tensor}. Unlike conventional trichromatic (RGB) images that integrate photon irradiance across three broad color channels, a hyperspectral image (HSI) records radiance across tens to hundreds of narrow, contiguous spectral bands spanning the visible and infrared spectra, facilitating non-invasive material identification and biochemical classification.

However, capturing hyperspectral data introduces severe physical trade-offs: splitting photon energy across narrow spectral bins sharply reduces photon counts at each detector pixel, yielding low signal-to-noise ratios (SNR), photon-counting Poisson shot noise, electronic Gaussian noise, sensor striping, and unobserved pixels \citep{bioucas2012hyperspectral, arad2016sparse}. Recovering an underlying clean tensor $\mathcal{Y} \in \mathbb{R}^{H \times W \times B}$ from a degraded observation $\mathcal{Y}^\text{noisy} \in \mathbb{R}^{H \times W \times B}$ constitutes an ill-posed inverse problem.

\subsection{Classical Multilinear Decompositions versus Deep Priors}
To regularize this inverse problem, classical signal processing leverages low-rank multilinear tensor models \citep{kolda2009tensor}. Exploiting spatial self-similarity across $(h, w)$ and spectral continuity across bands $b \in \{1, \dots, B\}$, the tensor $\mathcal{Y}$ is approximated via Canonical Polyadic (CP/PARAFAC) \citep{harshman1970foundations} or Tucker decompositions \citep{tucker1966some}. Standard algorithms (e.g., ALS) enforce low multilinear rank. Nevertheless, classical methods suffer from two key limitations: (i) \textbf{Rank Fragility:} Intrinsic tensor rank determination is NP-hard \citep{kolda2009tensor, heckel2020deeptensor}; underestimating rank truncates subtle spectral peaks, while overestimating rank fits sensor noise; and (ii) \textbf{Absence of Structural Priors:} Multilinear factorizations operate as algebraic contractions without inductive biases for spatial translation equivariance or edge continuity unless coupled with hand-crafted regularizers (e.g., Total Variation), which require delicate per-dataset tuning.

Supervised neural networks achieve high restoration fidelity, but acquiring paired ground-truth HSI datasets is frequently infeasible in specialized domains. To overcome this, \citet{ulyanov2018deep} introduced the \textit{Deep Image Prior} (DIP), demonstrating that an untrained convolutional network acts as an implicit regularizer: when fitted to a degraded image from random latent noise, low-frequency structures emerge far faster than high-frequency noise. However, extending DIP directly to 3D cubes via 3D CNNs is computationally intractable, incurring cubic memory complexity ($O(HWB)$) and GPU out-of-memory bottlenecks.

\subsection{The DeepTensor Framework}
To resolve this dilemma, \citet{heckel2020deeptensor, heckel2022deeptensor_arxiv} proposed \textbf{DeepTensor}, bridging low-rank tensor factorizations with deep image priors. Rather than generating the 3D cube via a 3D CNN, DeepTensor factors the matricized tensor $Y_{(1)} \in \mathbb{R}^{HW \times B}$ into low-rank spatial and spectral matrix components:
\begin{equation}
\label{eq:deeptensor_cp_intro}
\hat{\mathcal{Y}} = \sum_{r=1}^R \mathbf{u}_r \circ \mathbf{v}_r \quad \Longleftrightarrow \quad \hat{Y}_{(1)} = \mathbf{U} \mathbf{V}^T, \quad \mathbf{U} \in \mathbb{R}^{HW \times R}, \; \mathbf{V} \in \mathbb{R}^{B \times R}
\end{equation}
The factor matrices are parameterized by independent, untrained deep neural networks: $\mathbf{U} = \text{vec}(f_U(\mathbf{z}_U; \boldsymbol{\theta}_U))$ and $\mathbf{V} = f_V(\mathbf{z}_V; \boldsymbol{\theta}_V)$, where $f_U$ is a 2D CNN accepting fixed spatial noise $\mathbf{z}_U$, and $f_V$ is a 1D CNN accepting fixed spectral noise $\mathbf{z}_V$. Parameters are fitted by minimizing $\ell_1$ loss: $\min_{\boldsymbol{\theta}_U, \boldsymbol{\theta}_V} \|\mathbf{U}(\boldsymbol{\theta}_U) \mathbf{V}(\boldsymbol{\theta}_V)^T - Y^\text{noisy}\|_1$. By decoupling spatial and spectral dimensions, DeepTensor replaces 3D convolutions with efficient 2D and 1D operations, retaining deep structural priors while slashing memory requirements.

\subsection{Problem Statement and Objectives of NeuroTensor}
Despite its elegance, the original DeepTensor formulation has critical constraints: manual fixed rank specification, lack of factor parsimony, failure on incomplete data, unexamined heteroscedastic noise sensitivity, lack of spatial uncertainty quantification, and unexamined Transformer/attention inductive biases.

To address these gaps, \textbf{NeuroTensor} investigates nine concrete research objectives: (1) baseline reproduction on ICVL ($174 \times 164 \times 260$); (2) self-supervised \textbf{Adaptive Rank Selection}; (3) joint \textbf{Factor Group-Sparsity (Group Lasso)}; (4) \textbf{Mask-Aware Tensor Completion} ($20\%, 50\%, 80\%$); (5) \textbf{Noise Robustness} under matched input PSNR; (6) \textbf{Ensemble Disagreement} for spatial uncertainty; (7) \textbf{Factor Energy and Component Ablation}; (8) \textbf{Hybrid CNN--Attention} factor generation; and (9) pure \textbf{Spectral Transformer with ALiBi} relative biases across four diverse multimodal inverse problems.

\subsection{Research Contributions}
\begin{itemize}[noitemsep, topsep=1pt]
    \item \textbf{Nine Systematic Extensions:} Spanning rank adaptivity, group sparsity, completion, noise sensitivity, uncertainty, factor interpretability, and spectral Transformer architectures.
    \item \textbf{Unified 5-Model Benchmark:} Parameter-matched comparison: (1) Baseline CNN (\texttt{skip1d}), (2) Hybrid Attention (\texttt{hybrid}), (3) Conv Control (\texttt{cnn\_control}), (4) Pointwise Control (\texttt{pointwise\_control}), and (5) Pure Transformer (\texttt{transformer}).
    \item \textbf{Cross-Modality Inductive Bias Audit:} Demonstrating that CNNs excel in continuous spectral HSI ($46.07\text{ dB}$ vs.\ $43.40\text{ dB}$), Transformers lead in photometric face denoising ($28.56\text{ dB}$), and Pointwise Controls excel in tomography ($30.62\text{ dB}$).
    \item \textbf{Inductive Bias Mismatch Discovery:} Showing self-attention lacks translational smoothness in continuous spectra, incurring a $-2.55\text{ dB}$ penalty relative to 1D CNNs.
    \item \textbf{Fully Audited Reproducibility:} All code, derivations, and telemetry verified with zero fabricated claims.
\end{itemize}

\subsection{Report Organization}
Section~\ref{sec:related_work} reviews multilinear algebra, DIP, HSI restoration, and Transformers. Section~\ref{sec:data_methods} details dataset specifications, DeepTensor formulations, 5 architectures, and Extensions 1--9. Section~\ref{sec:results_discussion} presents experimental results, 12 comparative tables, multi-panel visual comparisons, and negative results analysis. Section~\ref{sec:conclusion} concludes and outlines future directions. Appendices A, B, and C provide architectural specifications, mathematical derivations, and execution telemetry.
